\documentclass[12pt]{article}
\usepackage{graphicx} %package to manage images
\graphicspath{ {./figures/} }
\usepackage{caption}
\usepackage[font=scriptsize]{subcaption}
\captionsetup[figure]{labelsep=none}
\captionsetup[table]{labelsep=none}
\usepackage{bbm}
\usepackage{amsmath}
\usepackage{import}
\usepackage{array}
\usepackage{enumitem}
\usepackage{booktabs}     % no extra vertical space between paragraphs
\usepackage{indentfirst}
\usepackage{afterpage}
\usepackage{floatrow}
\usepackage{pdflscape}
\usepackage{soul}
\usepackage{float}
\usepackage{adjustbox}
\usepackage{longtable}
\usepackage{caption}
\usepackage{setspace}
\usepackage{afterpage}
\usepackage[margin=1in]{geometry}
\usepackage[round]{natbib}
\usepackage{hyperref}
\usepackage{titlesec}
\usepackage{threeparttable}
\usepackage{setspace}
\usepackage{float}
\floatstyle{plaintop}
\restylefloat{table}
\usepackage{rotating}
\usepackage{ragged2e}
\usepackage{authblk}


\begin{document}
% \RaggedRight

% \textbf{Social group disparities in neonatal mortality are much larger among rural private births than public births, especially those in UP and Bihar.}

% \begin{table}[H]
%     \centering
%     \begin{threeparttable}
%         \caption{: Rural public sample - neonatal mortality disparities in NFHS-5}
%         \label{tab:regression}

%         \scriptsize
%         \setlength{\tabcolsep}{9pt}
%         \renewcommand{\arraystretch}{1.5}

%         % {\input{tables/predictor_regression}}
%         {\begin{tabular}{lccccc}
\hline
Region & Forward Hindu NNM & Adivasi gap & Dalit gap & OBC gap & Muslim gap \\
\hline
UP and Bihar&34.2&6.5&-0.9&-5.1&-0.7\\
EAG states besides UP and Bihar&22.2&7.0*&7.3*&0.4&-2.1\\
Non EAG states&16.3&5.0&2.0&3.2&1.1\\
\hline
\end{tabular}
}

%         \begin{tablenotes}
%             \item * $p<0.5$, ** $p<0.01$ *** $p<0.001$ 
%         \end{tablenotes}
%     \end{threeparttable}
% \end{table}




% \begin{table}[H]
%     \centering
%     \begin{threeparttable}
%         \caption{: Rural private sample - neonatal mortality disparities in NFHS-5}
%         \label{tab:regression}

%         \scriptsize
%         \setlength{\tabcolsep}{9pt}
%         \renewcommand{\arraystretch}{1.5}

%         % {\input{tables/predictor_regression}}
%         {\begin{tabular}{lccccc}
\hline
Region & Forward Hindu NNM & Adivasi gap & Dalit gap & OBC gap & Muslim gap \\
\hline
UP and Bihar&33.8&--&40.5***&11.7*&20.8**\\
EAG states besides UP and Bihar&19.5&25.7**&19.2**&0.3&!4.5\\
Non EAG states&11.6&8.7*&8.7**&3.0&-0.3\\
\hline
\end{tabular}
}

%         \begin{tablenotes}
%             \item * $p<0.5$, ** $p<0.01$ *** $p<0.001$ 
%         \end{tablenotes}
%     \end{threeparttable}
% \end{table}

\begin{figure}
    \centering
    \includegraphics[width=0.5\linewidth]{figures/}
    \caption{Caption}
    \label{fig:placeholder}
\end{figure}


\textbf{Large social group disparities in excess private neonatal mortality in UP and Bihar but not other states.} 


\begin{table}[H]
    \centering
    \begin{threeparttable}
        \caption{: Rural sample - excess private neonatal mortality (after controlling for selection) in the NFHS-4 and 5 pooled.}
        \label{tab:regression}

        \scriptsize
        \setlength{\tabcolsep}{9pt}
        \renewcommand{\arraystretch}{2}

        % {\input{tables/predictor_regression}}
        {\begin{tabular}{lccccc}
\hline
Region & Adivasi & Dalit & OBC & Forward Hindu & Muslim \\
\hline
UP and Bihar&--&39.7***&14.9***&6.2&28.9***\\
EAG states besides UP and Bihar&11.1&7.8&-3.5&3.4&!-6.4\\
Non EAG states&-0.9&0.5&-4.9*&-4.0&-2.7\\
\hline
\end{tabular}
}

        \begin{tablenotes}
            \item 
            \item - Each cell reports $\beta$ from a separate region-by-social-group regression of neonatal mortality on private facility birth:
            \[
            NNM_i = \alpha + \beta Private_i + X_i'\gamma + \delta_d + \alpha_r +\varepsilon_i,
            \]
            where $\delta_d$ are district fixed effects, and $\alpha_r$ are survey round fixed effects. Controls include prior neonatal death, projected maternal prepregnancy underweight, child sex, maternal age under 20, birth order, multiples, wealth quintile, and maternal illiteracy.
            \item
            \item - The sample in each regression is restricted to public and private facility births only.
            \item 
            \item - *, **, and *** indicate significance at the 10\%, 5\%, and 1\% levels, respectively. Standard errors are clustered by PSU.
            \item 
            \item - "--" indicates that the regression has fewer than 300 private births or fewer than 10 private neonatal deaths.
            \item - "!" indicates that the regression is flagged for having 300--500 private births or 10--20 private neonatal deaths.

        \end{tablenotes}
    \end{threeparttable}
\end{table}



\begin{itemize}
    \item If I do this table without controls and using care at birth indicators (early initiation of breastfeeding, skin to skin, prelacteal feeding), for all regions and all social groups the coefficient is large, negative, and *** significant. Those variables, at least in the NFHS, don't seem to be explaining why the excess private NNM is worse for disadvantaged social groups.
\end{itemize}








\end{document}